%----------------------------------------------------------------------------------------
%	Hamiltonian Flow Control (HFC) -- journal manuscript
%
%	Build:  latexmk main.tex   (XeLaTeX, see .latexmkrc; bibliography via biblatex + biber)
%	Notes:  notes/discussion.md (decisions, notation, open questions)
%	        notes/outline.md    (writing plan and progress)
%----------------------------------------------------------------------------------------

%----------------------------------------------------------------------------------------
%	PACKAGES AND DOCUMENT CONFIGURATIONS
%----------------------------------------------------------------------------------------

\documentclass[11pt,a4paper]{article}

\usepackage{amsmath}
\usepackage{amsthm}
\usepackage{mathtools}
\usepackage{fontspec}
\usepackage{xeCJK} % Only for Chinese drafting notes (\TODO); drop for submission
\usepackage{unicode-math}
\usepackage{geometry}
\usepackage{microtype}
\usepackage{graphicx}
\usepackage{subcaption}
\usepackage{booktabs}
\usepackage{multirow}
\usepackage{xcolor}
% Citation style: numbered in order of first citation, as in Nonlinear Dynamics.
% Cite with \citet (textual) and \citep (parenthetical): natbib=true provides them here, and
% the natbib-based Springer Nature template understands them too, so the text never changes.
\usepackage[
    backend=biber,
    style=numeric-comp,
    sorting=none,
    natbib=true,
    maxbibnames=99,
    giveninits=true,
    url=false,
    isbn=false,
]{biblatex}
\usepackage{hyperref}
\usepackage[capitalise,nameinlink]{cleveref}

%----------------------------------------------------------------------------------------
%	FONT SETTINGS
%----------------------------------------------------------------------------------------

% The fonts/ submodule stores its files with Git LFS. If only the LFS pointer files are
% checked out (git-lfs not installed), fall back to the identical New Computer Modern
% fonts shipped with TeX Live and to Noto CJK (same design as Source Han).
\IfFontExistsTF{fonts/NewCM10-Regular.otf}
    {\newcommand*{\fontdir}{fonts/}}
    {\newcommand*{\fontdir}{}}

\setmainfont{NewComputerModern}[
    Path=\fontdir,
    Extension=.otf,
    UprightFont=NewCM10-Regular,
    BoldFont=NewCM10-Bold,
    ItalicFont=NewCM10-Italic,
    BoldItalicFont=NewCM10-BoldItalic,
    FontFace={l}{n}{Font=NewCM10-Book.otf},
    FontFace={l}{it}{Font=NewCM10-BookItalic.otf},
]

\setsansfont{NewComputerModernSans}[
    Path=\fontdir,
    Extension=.otf,
    UprightFont=NewCMSans10-Regular,
    BoldFont=NewCMSans10-Bold,
    ItalicFont=NewCMSans10-Oblique,
    BoldItalicFont=NewCMSans10-BoldOblique,
    FontFace={l}{n}{Font=NewCMSans10-Book.otf},
    FontFace={l}{it}{Font=NewCMSans10-BookOblique.otf},
]

\setmathfont{NewCMMath-Regular.otf}[Path=\fontdir]

\IfFontExistsTF{fonts/SourceHanSerifSC-Regular.otf}{
    \setCJKmainfont{SourceHanSerifSC-Regular.otf}[
        Path=fonts/,
        BoldFont=SourceHanSerifSC-Bold.otf,
        ItalicFont=ukai.ttc,
        BoldItalicFont=ukai.ttc,
        FontFace={l}{n}{Font=SourceHanSerifSC-Light.otf},
    ]
    \setCJKsansfont{SourceHanSansSC-Regular.otf}[
        Path=fonts/,
        BoldFont=SourceHanSansSC-Bold.otf,
        FontFace={l}{n}{Font=SourceHanSansSC-Light.otf},
    ]
    \setCJKmonofont{SourceHanSansSC-Regular.otf}[Path=fonts/]
}{
    \setCJKmainfont{Noto Serif CJK SC}
    \setCJKsansfont{Noto Sans CJK SC}
    \setCJKmonofont{Noto Sans CJK SC}
}

%----------------------------------------------------------------------------------------
%	DOCUMENT SETTINGS
%----------------------------------------------------------------------------------------

\geometry{a4paper,margin=2.5cm}

\graphicspath{{figures/}}

\hypersetup{
    colorlinks=true,
    linkcolor=blue!50!black,
    citecolor=blue!50!black,
    urlcolor=blue!50!black,
    pdftitle={Hamiltonian Flow Control},
    pdfauthor={Houchen Li},
}

% Theorem-like environments, each numbered on its own (Theorem 1, Lemma 1, ...).
% Defined after cleveref so that \cref picks up the environment names.
\theoremstyle{plain}
\newtheorem{theorem}{Theorem}
\newtheorem{proposition}{Proposition}
\newtheorem{lemma}{Lemma}
\newtheorem{corollary}{Corollary}
\theoremstyle{definition}
\newtheorem{definition}{Definition}
\newtheorem{assumption}{Assumption}
\newtheorem{problem}{Problem}
\newtheorem{example}{Example}
\theoremstyle{remark}
\newtheorem{remark}{Remark}

\crefname{assumption}{Assumption}{Assumptions}
\crefname{problem}{Problem}{Problems}

%----------------------------------------------------------------------------------------
%	NOTATION MACROS
%	Single source of truth for notation; keep in sync with the symbol table in
%	notes/discussion.md.
%----------------------------------------------------------------------------------------

% Name of the method (session 18: Hamiltonian flow control). Change it here only; the title in
% main.tex and the pdftitle repeat it in full.
\newcommand{\method}{HFC}
\newcommand{\methodlong}{Hamiltonian flow control}

% Drafting placeholder for content still to be supplied or confirmed.
% Every occurrence must be resolved before submission.
\newcommand{\TODO}[1]{\textcolor{red}{[\,#1\,]}}

% Vectors are bold upright, as in the HFC derivations; scalars are italic.
\newcommand{\bx}{\mathbf{x}}
\newcommand{\bu}{\mathbf{u}}
\newcommand{\bv}{\mathbf{v}}
\newcommand{\bxref}{\mathbf{x}_{\mathrm{ref}}}

% Number sets
\newcommand{\R}{\mathbb{R}}
\newcommand{\N}{\mathbb{N}}

% Transpose: A^\T
\newcommand{\T}{\top}

% Operators
\DeclareMathOperator{\diag}{diag}
\DeclareMathOperator{\tr}{tr}

% Delimiters: \abs{x}, \norm{x}, \inner{x}{y}; starred forms auto-size
\DeclarePairedDelimiter{\abs}{\lvert}{\rvert}
\DeclarePairedDelimiter{\norm}{\lVert}{\rVert}
\DeclarePairedDelimiterX{\inner}[2]{\langle}{\rangle}{#1,#2}

% Upright differential: \int f(x) \dd x
\newcommand{\dd}{\mathop{}\!\mathrm{d}}

% Bold Greek vectors (unicode-math: \symbf instead of \boldsymbol)
\newcommand{\bxi}{\symbf{\xi}}
\newcommand{\bdelta}{\symbf{\delta}}
\newcommand{\blambda}{\symbf{\lambda}}


\addbibresource{ref/refs.bib}

% Allow looser lines in the bibliography rather than overfull boxes from long titles and DOIs.
\appto{\bibsetup}{\emergencystretch=1em}

%----------------------------------------------------------------------------------------
%	DOCUMENT INFORMATION
%----------------------------------------------------------------------------------------

\title{Hamiltonian Flow Control:\\[0.5ex]
    \large An Energy-Descent Control Law for Nonlinear Systems}
\author{Houchen Li\\[0.5ex]
    \normalsize\TODO{单位；通讯作者邮箱}}
\date{}

\begin{document}

\maketitle

%----------------------------------------------------------------------------------------
%	ABSTRACT
%----------------------------------------------------------------------------------------

\begin{abstract}
\TODO{摘要在全文完成后撰写：问题 — 方法 — 理论结果 — 三个算例与实验 — 结论；字数按目标期刊要求}

\medskip
\noindent\textit{Keywords:} \TODO{4--6 个关键词}
\end{abstract}


%----------------------------------------------------------------------------------------
%	SECTION 1: INTRODUCTION
%----------------------------------------------------------------------------------------

\section{Introduction}\label{sec:introduction}

% Paragraph plan: notes/outline.md, Section 1. The comments mark the role of each paragraph.

% (1) Background: industrial practice and the trade-off between cost and model use.
Feedback control has developed rapidly over the past decades, yet the algorithms deployed at
scale in industrial production remain few.
Proportional--integral--derivative (PID) control, the linear quadratic regulator (LQR), and
model predictive control (MPC) are the most representative of them: in a survey conducted by
the IFAC Industry Committee, PID control received the highest rating for current industrial
impact and MPC the highest rating for expected future impact~\citep{samad2020industry}.
These algorithms occupy different positions in a trade-off between computational cost and
the ability to exploit a model of the plant, that is, to anticipate its motion and to respect
its constraints.
Many emerging applications, such as agile robots, automated vehicles, and aerial vehicles,
demand both at once: their dynamics are nonlinear and often underactuated, and their
controllers must run at high rates on embedded processors with limited computing resources.

% (2) PID: model-free and cheap, but reactive and sensitive to its error history.
PID control is the representative of model-free feedback.
Its control signal combines the present error, the accumulated past error, and the rate of
change of the error, its tuning requires no model of the plant, and it remains the workhorse
of industrial feedback~\citep{astrom2001future}.
Its computation amounts to a few arithmetic operations per sample, so that it runs comfortably
on microcontrollers.
These advantages come with limitations.
Because the integral term stores the entire error history, past errors keep acting on the
plant long after they have occurred: transient disturbances and measurement noise that enter
the integrator are never forgotten, and while the actuator saturates, the integrator keeps
accumulating error that is released only after the actuator desaturates.
The result can be large overshoot and sustained oscillation, a phenomenon known as integrator
windup~\citep{astrom1989windup,astrom2008feedback}.
The derivative term, in turn, amplifies high-frequency measurement noise and must be
filtered~\citep{astrom2008feedback}.
Moreover, PID control is essentially reactive.
Its only anticipation is the linear extrapolation of the error implied by the derivative
term~\citep{astrom2008feedback}, so it can anticipate neither the nonlinear motion of the plant
nor upcoming changes of the reference, and its tracking performance deteriorates when the
reference varies rapidly.
Finally, although PID control appears conceptually simple, tuning it to meet several
conflicting objectives is hard in practice~\citep{ang2005pid}, and the difficulty is
compounded in strongly coupled multivariable and underactuated systems, where the loops are
designed channel by channel.

% (3) LQR: model-based and optimal, but Riccati equations, no constraints, local only.
The LQR is the prototype of model-based optimal
feedback~\citep{kalman1960contributions,anderson1990optimal}.
For a linear model and a quadratic cost, the optimal control is a linear state feedback whose
gain follows from a Riccati equation.
For linear time-invariant regulation, this equation is solved once, offline.
For nonlinear plants that operate over wide ranges or track time-varying references, however,
the model must be linearized along the operating trajectory and the Riccati equation solved
again online, or the gains must be scheduled~\citep{rugh2000research}, which raises the
computational burden markedly above that of PID control.
Standard LQR cannot enforce input or state constraints: saturating its output forfeits its
optimality and stability guarantees, and incorporating the constraints turns the control law
into a piecewise affine function of the state that must be computed by multiparametric
programming~\citep{bemporad2002explicit}.
LQR has further limitations.
When it is applied to a nonlinear plant through linearization, stability is guaranteed only in
a neighborhood of the operating point, so that large-amplitude maneuvers such as the swing-up
of a pendulum lie outside its scope.
The weighting matrices of the cost are only indirectly related to time-domain specifications
and are usually tuned by trial and error.
And the celebrated robustness margins of LQR~\citep{safonov1977gain} hold only under full-state
feedback: no margins are guaranteed once a state observer is inserted into the
loop~\citep{doyle1978guaranteed}.

% (4) MPC: constraints and nonlinear models, but an optimization problem per sample.
MPC overcomes several of these limitations by solving, at every sampling instant, a
finite-horizon optimal control problem for the current state and applying the first element of
the optimal input sequence~\citep{mayne2000constrained}.
It handles multivariable and nonlinear models, incorporates input and state constraints
explicitly, and has become the standard advanced control technology in the process
industries~\citep{qin2003survey}.
Its price is computational: a quadratic program, or a nonlinear program for nonlinear models,
must be solved within each sampling period.
Efficient solvers exploit the sparse, stage-wise structure of these problems through structured
or sparse matrix factorizations~\citep{rao1998application,wang2010fast,stellato2020osqp}, and
real-time schemes either perform only a single Newton-type iteration per
sample~\citep{diehl2005realtime,verschueren2022acados} or, as continuation methods do, update
the input sequence by integrating a differential equation that drives the residual of the
optimality conditions to zero~\citep{ohtsuka2004continuation}; nevertheless, the computational
and memory demands remain far above those of LQR.
Explicit MPC moves the optimization offline, but the number of regions of the resulting
piecewise affine law grows rapidly with the size of the problem~\citep{bemporad2002explicit}.
MPC therefore remains difficult to deploy on low-cost embedded platforms with fast sampling.
In addition, its closed-loop stability rests on suitably designed terminal costs and
constraints~\citep{mayne2000constrained}.

% (5) iLQR: nonlinear extension of LQR, at a higher computational cost.
The iterative linear quadratic regulator (iLQR)~\citep{li2004iterative} extends LQR to
nonlinear dynamics.
It is a variant of differential dynamic programming~\citep{mayne1966second} that neglects the
second-order derivatives of the dynamics: each iteration linearizes the dynamics around the
current trajectory, solves the resulting time-varying LQR problem by a backward Riccati-type
recursion, and updates the trajectory by a forward rollout.
Applied in a receding-horizon fashion, iLQR-type methods have synthesized complex behaviors
online, although for a humanoid robot the computation still ran about seven times slower than
real time on a standard PC~\citep{tassa2012synthesis}; constrained variants have been adopted
for motion planning in autonomous driving~\citep{chen2017constrained,chen2019autonomous}.
Every control update requires repeated forward--backward passes over the horizon, so that the
computational cost exceeds that of LQR; constraints are not handled natively and call for
dedicated extensions~\citep{tassa2014control,chen2019autonomous}; and, being a local method,
iLQR may converge to poor local solutions depending on its initialization.

% (6) Structured nonlinear designs: rigorous but system-specific.
Beyond these industrial workhorses, nonlinear control theory offers constructive designs with
rigorous guarantees, including feedback linearization~\citep{isidori1995nonlinear},
backstepping~\citep{krstic1995nonlinear}, sliding-mode control~\citep{utkin1977variable}, and
passivity-based control~\citep{ortega2002interconnection}; see \citet{khalil2002nonlinear} for
an overview.
These designs exploit structural properties of the model, for instance strict-feedback forms in
backstepping or the solvability of matching partial differential equations in the energy
shaping of underactuated mechanical systems~\citep{ortega2002stabilization}, and they are
largely system-specific.
Correspondingly, industrial surveys rate the current impact of nonlinear control far below that
of PID control and MPC~\citep{samad2020industry}.

% (7) Research gap.
The picture that emerges is one of two extremes.
PID control requires negligible computation but neither uses a model nor anticipates the motion
of the plant, whereas LQR, iLQR, and MPC exploit a model at a computational cost that grows from
solving Riccati equations to solving an optimization problem at every sampling instant.
What is missing is a controller that exploits a nonlinear model to anticipate the motion of the
plant, acts on every state component, including the unactuated states of underactuated systems,
keeps the per-step computation low enough for low-cost embedded hardware, and comes with
convergence guarantees.
\TODO{研究空白；“单步计算量低”需在实验中以实测耗时支撑（笔记 R11）}

% (8) Speed-gradient method: the starting point of HFC.
The speed-gradient (SG) method introduced by Fradkov~\citep{fradkov1979speed,fradkov1999nonlinear}
points in this direction.
Consider a plant $\dot\bx = \mathbf{f}(\bx,\bu,t)$ and a goal function $Q(\bx,t) \ge 0$ whose
vanishing encodes the control objective.
The SG method computes the rate of change, or \emph{speed}, of $Q$ along the trajectories of the
plant, $\omega(\bx,\bu,t) = \partial Q(\bx,t)/\partial t + \nabla_{\bx} Q(\bx,t)^\T
\mathbf{f}(\bx,\bu,t)$, and changes the input against the gradient of this speed with respect to
the input: in the differential form $\dot\bu = -K \nabla_{\bu}\omega(\bx,\bu,t)$ or, for plants
affine in the input, in the finite form $\bu = -K \nabla_{\bu}\omega(\bx,\bu,t)$, where $K$ is a
positive definite gain matrix.
Under a convexity condition on $\omega$ with respect to $\bu$ and an attainability condition, which
requires an input that renders $\omega$ sufficiently negative, the goal $Q \to 0$ is achieved,
with the goal function serving as a Lyapunov function~\citep{fradkov1999nonlinear}.
The method is constructive and model-based; its computation per step reduces to the evaluation
of a gradient, with neither a Riccati equation nor an online optimization; and it applies to
nonlinear, non-affine, and time-varying systems.
It has accordingly served as a unifying framework for adaptive control~\citep{fradkov1979speed},
energy-based control of oscillations~\citep{fradkov1996swinging}, stabilization of invariant
sets~\citep{shiriaev2000stabilization}, and control of chaos~\citep{fradkov2007cybernetical}.
Its limitations stem from its reliance on the instantaneous speed of the goal function.
First, $\nabla_{\bu}\omega$ receives no contribution from the state components whose rates do not
depend on the input.
In mechanical and kinematic models, the input acts directly only on the highest-order states,
such as velocities, so that errors in the low-order states, such as positions, drop out of the
control law.
The designer must then construct an auxiliary goal function, typically energy-like, whose speed
depends explicitly on the input; since the attainability condition essentially requires this
function to be a control Lyapunov function, the construction is difficult for general
underactuated systems.
Second, the resulting control is myopic: it decreases $Q$ at the current instant without
anticipating the future motion of the plant.
Third, achieving $Q \to 0$ does not in general imply convergence of the full state.
Energy-based swing-up laws, including those obtained by the SG
method~\citep{fradkov1996swinging}, regulate the energy of a pendulum to that of its upright
equilibrium and thereby steer the state toward the corresponding homoclinic orbit rather than to
the equilibrium itself, so that a separate local stabilizer must be switched in near the upright
position~\citep{astrom2000swinging,shiriaev2000global}.

% (9) Flow-based cost of Lan and Cvitanovic: the second ingredient of HFC.
A different line of work, originating in the computation of unstable periodic orbits, shows how a
cost can be built from the flow itself.
\citet{lan2004variational} proposed a variational method in which a trial loop in phase space is
deformed toward a periodic orbit.
Their cost measures, along the loop, the mismatch $\lambda\bv - \tilde{\bv}$ between the velocity
field $\bv$ of the flow, rescaled by a local time-scaling factor $\lambda$, and the tangent
$\tilde{\bv}$ of the loop; in discretized form, it compares the image of each loop point under
the flow over one time step with the next loop point.
The loop evolves under a ``Newton descent'' equation, an infinitesimal-step version of the damped
Newton--Raphson iteration, along which this mismatch decays exponentially in a fictitious time,
so that the cost decreases monotonically.
For systems with higher time derivatives, the method matches the highest derivatives instead.
The idea of measuring the cost through the phase-space flow, rather than through a hand-crafted
function of the state, is the one we adopt in this paper.

% (10) Controllers whose input evolves as a flow: closest related work.
Letting the control input evolve as the solution of a differential equation that drives a
residual to zero is not new in control either.
Besides the continuation methods for receding-horizon control mentioned
above~\citep{ohtsuka2004continuation}, Wardi and co-workers have developed an output-tracking
controller based on a fluid-flow version of the Newton--Raphson method: an integrator with variable
gain that drives to zero the difference between a predicted future output and the reference
signal~\citep{wardi2017performance,wardi2024tracking}.
Implementable by simple algorithms, it has been applied to autonomous
vehicles~\citep{shivam2019tracking}, to aggressive quadrotor flight, where it outperformed the native
tracking controller of the PX4 autopilot~\citep{morales2024newton}, and to computationally
constrained aerial vehicles, where it matched or exceeded nonlinear MPC at a substantially lower
computation time~\citep{morales2026lightweight}; stability conditions are available for linear
systems and for a class of differentially flat systems~\citep{wardi2024tracking,niu2025stability}.
\method{} shares with this controller the idea of steering a predicted state toward the reference by
a flow of the input, but it differs in two respects.
First, the Newton--Raphson flow solves a square system of equations, with as many tracked outputs as
inputs, and its update requires the inverse of the Jacobian of the predicted output with respect to
the input.
\method{} instead weighs the predicted error of every state component and drives the input along a
gradient scaled by diagonal coefficients, so that no matrix is inverted and the states outside the
tracked outputs, such as the pendulum angles of an underactuated system, enter the control law
directly.
Second, in its basic form the Newton--Raphson flow predicts the output by integrating the plant over
the preview horizon with the input held constant~\citep{wardi2024tracking}, whereas \method{}
evaluates a Taylor polynomial of the flow in closed form, whose higher-order terms already make every
state component depend on the input.

% (11) The proposed approach.
In this paper, we propose \methodlong{} (\method{}), which builds on the SG method and on the
flow-based cost of \citet{lan2004variational}.
The reference is a sequence of state samples without time stamps, whose evolution times are close
to the update interval $h$ of the controller but not exactly uniform; at every round, the controller
takes the $m$th next sample as target, which requires no projection onto the reference, and aims to
reach it after the preview time $\tau = mh$.
Writing $\bv_\tau$ for the mean velocity of a look-ahead of the plant over the preview time, the cost
\[
    F(\bx,\bu,t) = \norm[\big]{\Gamma\bigl(\tau\bv_\tau(\bx,\bu) - (\bxref - \bx)\bigr)}^2
\]
compares the flow of the plant with the mean velocity that would carry the current state to the
target sample, the matrix $\Gamma$ rendering the cost dimensionless.
This is the counterpart of the mismatch $\lambda\bv - \tilde{\bv}$ of \citet{lan2004variational}, the
secant toward the target sample playing the role of the loop tangent and the preview time that of
the time-scaling factor: the controller chases a flow rather than a geometric position.
The control input evolves along the negative gradient of $F$ with respect to $\bu$, scaled by
dimension-balancing coefficients, and the controller follows this gradient flow by one explicit
update per round, without any online optimization.
To first order in $\tau$, this control law coincides, up to the gain matrix, with the SG law for the
quadratic goal function $Q = \frac{1}{2}\norm{\Gamma(\bx - \bxref)}^2$ (\cref{prop:sg}), and a
look-ahead of first order would therefore inherit the first limitation of SG: the low-order states
would not contribute to the gradient.
\method{} removes this limitation by taking as look-ahead the Taylor polynomial of the flow of the
plant, with the input held over the preview time, up to the order at which the input reaches every
state component: the low-order states are predicted through the time derivatives of the high-order
ones, much as \citet{lan2004variational} match the highest derivatives of systems of higher order,
so that every state component contributes to the control update.
As a result, and unlike SG, \method{} needs no hand-crafted energy-like goal function: a quadratic
cost on the state errors suffices.
Expanding the square, $F/2 = Q + \tau\mathcal{H}$, where $\mathcal{H}$ has the form of the Hamiltonian
of Pontryagin's principle~\citep{pontryagin1962mathematical}, with the gradient of $Q$ in the role of
the costate, the mean velocity of the look-ahead in that of the velocity, and the kinetic energy of
this velocity as running cost (\cref{rem:hamiltonian}).
The input therefore descends a control Hamiltonian continuously, rather than minimizing it along the
solution of an adjoint equation, which motivates the name of the method.
We call $F$ the flow cost; as the sum of the potential $Q$ and $\tau$ times this Hamiltonian, it is an
energy-like function, and \method{} is an energy-descent control law.
The Hamiltonian is not meant in the sense of Hamiltonian mechanics: the closed loop is dissipative,
and no symplectic structure is claimed.
The controller uses the target samples only: neither their timing, nor values between them, nor
the input that generates them is required.

% (12) Contributions.
The main contributions of this paper are as follows.
\begin{itemize}
    \item We propose \method{}, which drives the control input along the gradient of a cost that
    compares the flows of the plant and of the reference over a preview time, in the spirit of the
    variational method of \citet{lan2004variational}; we show that this cost is the weighted state
    error plus the preview time times a control Hamiltonian of Pontryagin type, and that the control
    law contains the SG law as its first-order approximation in the preview time.
    \item We introduce a Taylor look-ahead, through which the low-order states, which the input
    reaches only through the higher-order ones, contribute to the control update, so that no
    hand-crafted energy-like goal function is needed; the look-ahead is affine in the input, so
    that the cost is quadratic in it.
    \item We show that on integrator chains the closed-loop eigenvalues of \method{} are $1/\tau$
    times the roots of a polynomial fixed by the weights, which yields exact conditions for
    stability, met by equal weights up to chains of four levels, and for the update interval; we
    show that the look-ahead toward a target sample causes a time delay equal to the mismatch
    between the preview time and the time to the sample, with no offset up to terms of third order
    in the preview time or of second order in the rate of the reference input, that taking the $m$th
    next sample with $\tau = mh$ removes the delay without any projection onto the reference, and we
    bound the resulting tracking error.
    \TODO{一般非线性系统的收敛性定理（类型待定，见 notes/convergence.md）}
    \item We derive \method{} laws for three benchmark problems: the swing-up and stabilization of a
    double inverted pendulum on a cart~\citep{graichen2007swingup} by a single control law
    without switching, trajectory tracking for a kinematic bicycle model, for which we prove local
    exponential stability of straight references at every speed in continuous time and quantify the
    speed range of the sampled-data implementation, and trajectory tracking for a
    quadrotor whose attitude evolves on $\mathrm{SO}(3)$~\citep{lee2010geometric}.
    \item We validate \method{} in simulations and physical experiments on all three platforms.
    \TODO{对比基线与评价指标待定；建议包括 SG 与 Newton–Raphson 流方法（笔记 R6）}
\end{itemize}

% (13) Organization. Draft -- revisit once Sections 2--4 are settled.
The remainder of this paper is organized as follows.
\Cref{sec:formulation} recalls the required geometric preliminaries, introduces the flow cost and
its Taylor look-ahead, defines the control law and its sampled-data implementation, relates the
cost to a control Hamiltonian and the law to the SG method, analyzes the convergence and the
tracking error caused by the look-ahead, and treats non-uniform sample intervals; it then derives
the control laws for the three benchmark problems.
\Cref{sec:experiments} reports simulation and experimental results for these systems, and
\cref{sec:conclusion} concludes the paper and outlines possible extensions.

% (14) Notation.
\paragraph{Notation.}
\TODO{符号说明：待 Formulation 定稿后补充}


%----------------------------------------------------------------------------------------
%	SECTION 2: FORMULATION
%----------------------------------------------------------------------------------------

\section{Formulation}\label{sec:formulation}

% Plan and required inputs: notes/outline.md, Section 2. Since session 18 the look-ahead is the
% Taylor polynomial of the flow, as in kinematics_flow_controller.hpp; the staged (symplectic
% Euler) look-ahead of sessions 8-17 survives only in Remark rem:staged. Subsections 2.1, 2.9,
% 2.10 and 2.12 still carry only their plan.

\subsection{Geometric preliminaries}\label{sec:preliminaries}

% Only what the general theory and the three examples need: smooth manifolds, tangent and
% cotangent bundles, Lie derivatives along vector fields (Taylor look-ahead), the Lie groups SO(3)
% and SE(3) with the hat/vee maps (Taylor look-ahead on SO(3), Magnus-type expansion).
\TODO{流形、切丛与余切丛、李导数、李群 SO(3)/SE(3) 及 hat/vee 映射等预备知识，篇幅从简}

\subsection{Problem setting}\label{sec:problem}

We consider plants of the form
\begin{equation}\label{eq:plant}
    \dot\bx = \bv(\bx,\bu), \qquad \bx \in \R^n, \quad \bu \in \mathcal{U} \subseteq \R^m,
\end{equation}
with a smooth vector field $\bv$.
\TODO{本节暂在局部坐标中表述；流形（含 SO(3)）上的内蕴表述待 2.1 节完成后统一}
The state is partitioned into $L$ levels of increasing differential order,
\begin{equation}\label{eq:levels}
    \bx = (\bx_1,\dots,\bx_L), \qquad \bv = (\bv_1,\dots,\bv_L), \qquad \bx_\ell \in \R^{n_\ell},
\end{equation}
so that positions precede velocities, attitudes precede angular velocities, and so on; for $L = 2$
we also write $(\mathbf{q},\mathbf{p}) = (\bx_1,\bx_2)$.
The partition only orders the state components by the number of integrations that separate them
from the input; in particular, $\mathbf{q}$ and $\mathbf{p}$ need not be canonically conjugate.

The reference is given as a sequence of state samples $\bx_{r,0},\bx_{r,1},\dots$, as produced by a
planner, without time stamps.
The evolution times between consecutive samples are roughly stable, close to the update interval
$h$ of the controller, but not exactly uniform, and neither values between the samples nor the
input that generates them are assumed to be available.
At round $k$ of the controller, at time $t_k$, the target is the $m$th next sample,
\begin{equation}\label{eq:xref}
    \bxref(t) = \bx_{r,k+m}, \qquad t\in[t_k,t_{k+1}),
\end{equation}
which the controller aims to reach after the preview time $\tau = mh$, the time in which the
reference is expected to reach it; the integer $m \ge 2$ is a design parameter
(\cref{rem:target}), and taking a sample at a fixed index offset avoids any projection of the state
onto the reference.
The tracking error at round $k$ is the deviation of the state from the sample due at that round,
which is the target set $m$ rounds earlier,
\begin{equation}\label{eq:error}
    \mathbf{e}_k := \bx(t_k) - \bx_{r,k};
\end{equation}
we do not split it into along-track and cross-track components.
For the analysis only, we regard the samples as taken from a trajectory $\bx_r(t)$ of the plant,
$\dot\bx_r = \bv(\bx_r,\bu_r)$ for some input $\bu_r(t)$, with $\bx_{r,k} = \bx_r(t_k)$; the controller
never uses $\bx_r$ or $\bu_r$.
Set-point regulation is the special case of a constant reference.

\subsection{The flow-based cost}\label{sec:cost}

Given a look-ahead map $\Phi_\tau\colon \R^n\times\mathcal{U}\to\R^n$, which predicts the state reached
after the preview time under the input $\bu$, \method{} penalizes the mismatch between the
prediction and the target sample,
\begin{equation}\label{eq:cost}
    F(\bx,\bu,t) = \norm[\big]{\Gamma\bigl(\Phi_\tau(\bx,\bu) - \bxref(t)\bigr)}^2,
\end{equation}
where the invertible matrix $\Gamma$ renders every component dimensionless; we write
$W := \Gamma^\T\Gamma$.
With the mean velocity of the look-ahead, $\bv_\tau(\bx,\bu) := (\Phi_\tau(\bx,\bu) - \bx)/\tau$, the
cost reads
\begin{equation}\label{eq:cost-flow}
    F = \tau^2\,\norm[\Big]{\Gamma\Bigl(\bv_\tau(\bx,\bu) - \frac{\bxref - \bx}{\tau}\Bigr)}^2 ,
\end{equation}
so it compares the flow of the plant over the preview time with the mean velocity that would carry
the current state to the target sample in that time.
This is the counterpart of the cost $(\lambda\bv - \tilde{\bv})^2$ of \citet{lan2004variational}: in
the discretized form of their method, the image of a loop point under the flow over one time step is
compared with the next loop point, the loop tangent $\tilde{\bv}$ being the secant to that point; here
the look-ahead of the plant state is compared with the target sample, and the preview time plays the
role of the time-scaling factor $\lambda$.
We use the forward Euler look-ahead
\begin{equation}\label{eq:fe}
    \Phi_\tau^{\mathrm{FE}}(\bx,\bu) = \bx + \tau\bv(\bx,\bu)
\end{equation}
as a point of comparison below.
In practice $\Gamma$ is diagonal, $\Gamma = \diag(\sqrt{w_1},\dots,\sqrt{w_n})\,\bar\Gamma$, where
$\bar\Gamma$ divides each component by a characteristic scale built from the plant parameters and
$\tau$, and the weights $w_i \ge 0$ are pure numbers.

\subsection{Taylor look-ahead}\label{sec:lookahead}

With the forward Euler look-ahead, an input that drives only the higher levels cannot influence the
predicted lower levels: if $\partial_{\bu}\bv_\ell = 0$, the $\ell$th block row of
$\partial_{\bu}\Phi_\tau^{\mathrm{FE}}$ vanishes, and the error of level $\ell$ drops out of
$\nabla_{\bu}F$.
\method{} therefore predicts the lower levels through the time derivatives of the higher ones.
For a fixed input $\bu$, let $\mathfrak{L}_{\bu}$ denote the Lie derivative along $\bv(\cdot,\bu)$, so
that $\mathfrak{L}_{\bu}^j\bx$ is the $j$th time derivative of the state along the solutions of
\eqref{eq:plant} with the input held at $\bu$; thus $\mathfrak{L}_{\bu}\bx = \bv(\bx,\bu)$ and
$\mathfrak{L}_{\bu}^2\bx = \partial_{\bx}\bv(\bx,\bu)\,\bv(\bx,\bu)$.

\begin{definition}[Taylor look-ahead]\label{def:taylor}
The Taylor look-ahead of order $L$ is
\begin{equation}\label{eq:taylor}
    \hat\Phi_\tau(\bx,\bu) = \sum_{j=0}^{L}\frac{\tau^j}{j!}\,\mathbf{c}_j(\bx,\bu),
\end{equation}
where $\mathbf{c}_j = \mathfrak{L}_{\bu}^j\bx$ for $j \le 2$ and, for $j \ge 3$, $\mathbf{c}_j$ is the
first-order expansion of $\mathfrak{L}_{\bu}^j\bx$ in $\bu$ at $\bu = \mathbf{0}$.
\end{definition}

Up to order two, \eqref{eq:taylor} is the Taylor polynomial of the flow of \eqref{eq:plant} over the
preview time with the input held; from order three on, it keeps only the part of each term that is
affine in the input.
If the input drives only the highest level, $\partial_{\bu}\bv_\ell = 0$ for $\ell < L$, it first
appears in the derivatives of level $\ell$ at order $L - \ell + 1$, so that $L$ is the lowest order at
which it reaches every level.
On an integrator chain, $\hat\Phi_\tau$ is the exact flow.
If the plant is affine in the input, $\bv = \mathbf{f}(\bx) + \sum_i \mathbf{g}_i(\bx)u_i$, and the input
vector fields do not depend on the components they drive, $\partial_{\bx}\mathbf{g}_i\,\mathbf{g}_j =
\mathbf{0}$ for all $i$ and $j$, then $\mathbf{c}_1$ and $\mathbf{c}_2$ are affine in $\bu$ as well, so
that $\hat\Phi_\tau$ is affine in $\bu$ and $F$ is a quadratic function of $\bu$.
The three plants of this paper satisfy this condition; in the quadrotor, for instance, the thrust
acts along the attitude, which no input drives directly.
Expanding \eqref{eq:taylor} in $\tau$ gives
\begin{equation}\label{eq:taylor-expansion}
    \hat\Phi_\tau(\bx,\bu) = \bx + \tau\bv(\bx,\bu) + \tau^2\rho(\bx,\bu,\tau), \qquad
    \rho(\bx,\bu,0) = \tfrac12\,\partial_{\bx}\bv(\bx,\bu)\,\bv(\bx,\bu),
\end{equation}
where $\rho$ is a polynomial in $\tau$; its second-order term is what lets the input reach the levels
below the directly driven ones (\cref{prop:sg}).

\begin{remark}[Staged look-ahead]\label{rem:staged}
A cheaper look-ahead, which needs only the first derivatives of $\bv$, stages a single explicit step
in the order of the symplectic Euler scheme~\citep{hairer2006geometric}: every level is advanced over
the preview time with the already advanced higher levels,
$\bx_\ell^+ = \bx_\ell + \tau\bv_\ell(\bx_1,\dots,\bx_\ell,\bx_{\ell+1}^+,\dots,\bx_L^+,\bu)$, from
$\ell = L$ down to $\ell = 1$.
It also lets the input reach every level, but it holds the advanced higher levels over the whole
preview time.
On an integrator chain it gives $x_\ell^+ = \sum_{j\ge\ell}\tau^{j-\ell}x_j + \tau^{L-\ell+1}u$, with
all coefficients equal to one instead of $1/(j - \ell)!$, so that it overstates the reach of the input
into level $\ell$ by the factor $(L - \ell + 1)!$, and its second-order bias leaves a steady offset
(\cref{ex:offset}).
We use the Taylor look-ahead throughout.
\end{remark}

\subsection{Control law and implementation}\label{sec:extended}

The control input becomes a state of the controller and descends the cost:
\begin{equation}\label{eq:hfc}
    \dot\bx = \bv(\bx,\bu), \qquad \dot\bu = -D(\bx,\bu)\,\nabla_{\bu}F(\bx,\bu,t),
\end{equation}
where $F$ is built on $\hat\Phi_\tau$ and the dimension-balancing coefficients form the diagonal
matrix
\begin{equation}\label{eq:D}
    D = \diag(d_1,\dots,d_m), \qquad
    d_i = \Bigl(2\tau^3\,\norm[\big]{\bar\Gamma\,\partial_{u_i}\bv(\bx,\bu)}^2\Bigr)^{-1}.
\end{equation}
The coefficient $d_i$ is $1/\tau$ times the inverse curvature, with respect to $u_i$, of the
unweighted forward Euler cost $\norm{\bar\Gamma(\bx + \tau\bv - \bxref)}^2$.
It gives $\dot u_i$ the units of a rate of the input, and it scales the gradient so that a residual
on which $u_i$ acts directly, weighted with $w = 1$, decays at the rate $1/\tau$: the infinitesimal
counterpart of a full Newton step taken over the preview time.
The preview time therefore sets the loop gain, and the weights are pure numbers.

The controller runs in discrete time.
At round $k$, it measures the state $\bx(t_k)$ and updates the input by a forward Euler step of
\eqref{eq:hfc}, clamped into $\mathcal{U}$,
\begin{equation}\label{eq:update}
    \bu_{k+1} = \Pi_{\mathcal{U}}\Bigl(\bu_k - h\,D\bigl(\bx(t_k),\bu_k\bigr)\,
        \nabla_{\bu}F\bigl(\bx(t_k),\bu_k,t_k\bigr)\Bigr),
\end{equation}
where $\Pi_{\mathcal{U}}$ clamps every component into its bounds.
The plant moves under the input $\bu_k$, held over $[t_k,t_{k+1})$, and $\bu_{k+1}$ takes over at
$t_{k+1}$, which leaves one round for the computation.
We analyze this sampled-data loop below, and we simulate it with an accurate integration of the
plant.
The update may also be taken from a step of a higher-order scheme applied to \eqref{eq:hfc} from the
measured state, the stage at $t_k + c\,h$ then looking ahead over $\tau - c\,h$; the forward Euler
update is the cheapest, and it is the one we use in the experiments.
\TODO{约束的连续时间表述（投影）}

\subsection{Properties of the flow cost}\label{sec:properties}

\TODO{F 的非负性与零点集}

The flow cost has the structure of a control Hamiltonian, and it generalizes the SG method.
Let $Q(\bx,t) = \frac12\norm{\Gamma(\bx - \bxref(t))}^2$, whose gradient
$\nabla_{\bx}Q = W(\bx - \bxref)$ we denote by $\blambda$, and let
$\omega(\bx,\bu,t) = \partial_t Q(\bx,t) + \nabla_{\bx}Q(\bx,t)^\T\bv(\bx,\bu)$ be the speed of $Q$
along \eqref{eq:plant}, so that $\nabla_{\bu}\omega = (\partial_{\bu}\bv)^\T W(\bx - \bxref)$ is the
direction of the SG law for the goal function $Q$.

\begin{proposition}[Hamiltonian form and relation to the speed-gradient law]\label{prop:sg}
Let $\Gamma$ be fixed, write $K(\bx,\bu) = \frac12\norm{\Gamma\bv(\bx,\bu)}^2$, and define the control
Hamiltonian
\begin{equation}\label{eq:hamiltonian}
    \mathcal{H}(\bx,\bu,\blambda) = \blambda^\T\bv_\tau(\bx,\bu)
    + \frac{\tau}{2}\,\norm[\big]{\Gamma\bv_\tau(\bx,\bu)}^2 .
\end{equation}
\begin{enumerate}
    \item For every look-ahead, $F = 2Q + 2\tau\,\mathcal{H}(\bx,\bu,\nabla_{\bx}Q)$, and therefore
    $\nabla_{\bu}F = 2\tau\,\nabla_{\bu}\mathcal{H}$.
    \item For the forward Euler look-ahead \eqref{eq:fe}, $\mathcal{H} = \omega - \partial_t Q + \tau K$,
    and therefore
    \begin{equation}\label{eq:sg-fe}
        \nabla_{\bu}F = 2\tau\,\nabla_{\bu}\omega + 2\tau^2\,\nabla_{\bu}K.
    \end{equation}
    \item For the Taylor look-ahead \eqref{eq:taylor},
    \begin{equation}\label{eq:sg-taylor}
        \nabla_{\bu}F = 2\tau\,\nabla_{\bu}\omega
        + 2\tau^2\Bigl(\nabla_{\bu}K + \partial_{\bu}\rho(\bx,\bu,0)^\T W(\bx - \bxref)\Bigr)
        + O(\tau^3)
    \end{equation}
    uniformly on compact sets.
    Hence $\dot\bu = -2\tau D\bigl(\nabla_{\bu}\omega + O(\tau)\bigr)$: to first order in $\tau$,
    the \method{} law \eqref{eq:hfc} is the differential SG law with the gain matrix $2\tau D$.
    \item Let $W = \diag(W_1,\dots,W_L)$ be block diagonal, and let level $\ell$ satisfy
    $\partial_{\bu}\bv_\ell = 0$. Then the error $\mathbf{e}_\ell = \bx_\ell - \bx_{\mathrm{ref},\ell}$
    enters neither $\nabla_{\bu}\omega$ nor, for the forward Euler look-ahead, $\nabla_{\bu}F$.
    For the Taylor look-ahead, it enters $\nabla_{\bu}F$, up to terms of order $\tau^3$, through
    \begin{equation}\label{eq:sg-lower}
        \tau^2\Bigl(\sum_{k=1}^{L}\partial_{\bx_k}\bv_\ell\,\partial_{\bu}\bv_k\Bigr)^{\!\T}
        W_\ell\,\mathbf{e}_\ell .
    \end{equation}
\end{enumerate}
\end{proposition}

\begin{proof}
(a) Since $\Phi_\tau - \bxref = (\bx - \bxref) + \tau\bv_\tau$, expanding the square in
\eqref{eq:cost} gives
$F = \norm{\Gamma(\bx - \bxref)}^2 + 2\tau(\bx - \bxref)^\T W\bv_\tau + \tau^2\norm{\Gamma\bv_\tau}^2
= 2Q + 2\tau\,\mathcal{H}(\bx,\bu,\nabla_{\bx}Q)$, and $Q$ does not depend on $\bu$.

(b) For the forward Euler look-ahead, $\bv_\tau = \bv$ and
$(\bx - \bxref)^\T W\bv = \nabla_{\bx}Q^\T\bv = \omega - \partial_t Q$; neither $Q$ nor $\partial_t Q$
depends on $\bu$, which gives \eqref{eq:sg-fe}.

(c) With $\hat\Phi_\tau - \bxref = (\bx - \bxref) + \tau\bv + \tau^2\rho$,
\[
    \nabla_{\bu}F
    = 2\bigl(\tau\,\partial_{\bu}\bv + \tau^2\partial_{\bu}\rho\bigr)^\T W
      \bigl((\bx - \bxref) + \tau\bv + \tau^2\rho\bigr).
\]
The terms of order $\tau$ and $\tau^2$ are
\[
    2\tau(\partial_{\bu}\bv)^\T W(\bx - \bxref) = 2\tau\nabla_{\bu}\omega
    \quad\text{and}\quad
    2\tau^2\bigl((\partial_{\bu}\bv)^\T W\bv + (\partial_{\bu}\rho)^\T W(\bx - \bxref)\bigr),
\]
where $(\partial_{\bu}\bv)^\T W\bv = \nabla_{\bu}K$.
The remaining terms carry at least the factor $\tau^3$.
Since $\rho$ is a polynomial in $\tau$ whose coefficients are continuously differentiable in
$(\bx,\bu)$, replacing $\partial_{\bu}\rho(\bx,\bu,\tau)$ by $\partial_{\bu}\rho(\bx,\bu,0)$ changes the
$\tau^2$ term by $O(\tau^3)$.
All bounds are uniform on compact sets, where the functions involved are continuous.
Multiplying by $D$ gives the form of $\dot\bu$.

(d) Since $W$ is block diagonal, $\nabla_{\bu}\omega = \sum_k (\partial_{\bu}\bv_k)^\T
W_k\mathbf{e}_k$.
If $\partial_{\bu}\bv_\ell = 0$, the term with $k = \ell$ vanishes, so $\mathbf{e}_\ell$ does not appear
in $\nabla_{\bu}\omega$, nor, by \eqref{eq:sg-fe}, in $\nabla_{\bu}F$ for the forward Euler
look-ahead.
For the Taylor look-ahead, $\partial_{\bx_k}\bv_\ell$ does not depend on $\bu$, so that
\eqref{eq:taylor-expansion} gives $\partial_{\bu}\rho_\ell(\bx,\bu,0) =
\frac12\sum_k\partial_{\bx_k}\bv_\ell\,\partial_{\bu}\bv_k$, and the $\ell$th block of the
$\tau^2$ term in \eqref{eq:sg-taylor} is \eqref{eq:sg-lower}.
\end{proof}

\begin{remark}[Hamiltonian and energy-like structure]\label{rem:hamiltonian}
The function \eqref{eq:hamiltonian} has the form of the Hamiltonian of Pontryagin's
principle~\citep{pontryagin1962mathematical}, a running cost plus a costate times the velocity: the
gradient $\blambda = \nabla_{\bx}Q$ of the weighted state error plays the role of the costate, the
mean velocity $\bv_\tau$ of the look-ahead that of the velocity, and $\frac{\tau}{2}\norm{\Gamma
\bv_\tau}^2$, the kinetic energy of the mean velocity in the metric $W$, that of the running cost.
Pontryagin's principle requires an optimal input to minimize the Hamiltonian along the solution of
the adjoint equation.
By \cref{prop:sg}(a), \method{} instead lets the input descend $\mathcal{H}$ continuously, with the
gradient of $Q$ in place of the costate, so that $Q$ plays the role of a value function, and it needs
neither an adjoint equation nor a boundary-value problem; hence the name of the method.
Within a round the target is held, $\partial_t Q = 0$, and $F/2 = Q + \tau\mathcal{H}$ adds $\tau$
times this Hamiltonian to the potential $Q$; for the forward Euler look-ahead, $F/2 = Q + \tau\omega +
\tau^2 K$ is the sum of the potential, its rate along the flow, and a kinetic term.
The input thus descends an energy-like function of the state and of its flow.
The name does not refer to Hamiltonian mechanics: with the state frozen, $F$ decreases along
\eqref{eq:hfc} at the rate $(\nabla_{\bu}F)^\T D\,\nabla_{\bu}F$, so that the closed loop is
dissipative, and no symplectic structure is claimed.
\end{remark}

\begin{remark}\label{rem:sg-levels}
If level $\ell$ is separated from the input by more than one level, the sum in \eqref{eq:sg-lower}
vanishes and $\mathbf{e}_\ell$ enters $\nabla_{\bu}F$ at a higher order in $\tau$; on an integrator
chain, the error of the level that lies $j$ levels below the input enters at order $\tau^{j+1}$
(see the proof of \cref{prop:chain}).
The non-dimensionalization by $\bar\Gamma$, which scales the components of different levels by
different powers of $\tau$, makes these contributions comparable in size.
\end{remark}

\subsection{Convergence analysis}\label{sec:convergence}

\TODO{一般非线性系统的收敛性主定理，研究方案见 notes/convergence.md}

\paragraph{Integrator chains.}
The integrator chain is the prototype of the level structure \eqref{eq:levels}; on it, the Taylor
look-ahead is the exact flow, and the \method{} loop can be analyzed exactly.

\begin{proposition}[\method{} on an integrator chain]\label{prop:chain}
Consider the chain $\dot x_1 = x_2,\dots,\dot x_{L-1} = x_L$, $\dot x_L = u$ with scalar states and
input, the Taylor look-ahead \eqref{eq:taylor}, the coefficient \eqref{eq:D}, and
\[
    \Gamma = \ell_0^{-1}\diag\bigl(\sqrt{w_1},\ \sqrt{w_2}\,\tau,\ \dots,\ \sqrt{w_L}\,\tau^{L-1}\bigr)
\]
with a characteristic length $\ell_0$ and weights $w_k > 0$.
For set-point regulation, $\bx_r = 0$ and $u_r = 0$, the extended system \eqref{eq:hfc} is linear
and its characteristic polynomial is $\tau^{-(L+1)}P(\tau s)$, where
\begin{equation}\label{eq:chain-poly}
    P(\sigma) = \sigma^{L+1} + \sum_{k=1}^{L}\frac{w_k}{(L - k + 1)!}\,\sigma^{k-1}e_{L-k+1}(\sigma),
    \qquad e_n(\sigma) = \sum_{i=0}^{n}\frac{\sigma^i}{i!} .
\end{equation}
Consequently, the closed loop is exponentially stable if and only if $P$ is Hurwitz, a condition on
the weights alone, and every closed-loop eigenvalue is $1/\tau$ times a root of $P$.
\end{proposition}

\begin{proof}
Since the chain is linear, \eqref{eq:taylor} is its exact flow,
$x_k^+ = \sum_{j=k}^{L}\tau^{j-k}x_j/(j-k)! + \tau^{L-k+1}u/(L-k+1)!$.
In the scaled variables $z_j = \tau^{j-1}x_j$ and $\nu = \tau^L u$ this reads $\tau^{k-1}x_k^+ = Y_k$
with $Y_k = \sum_{j\ge k} z_j/(j-k)! + \nu/(L-k+1)!$, so that $F = \ell_0^{-2}\sum_{k} w_k Y_k^2$.
Here $\bar\Gamma\,\partial_u\bv = \tau^{L-1}/\ell_0$, so \eqref{eq:D} gives
$d = \ell_0^2/(2\tau^{2L+1})$ and
$\dot u = -d\,\partial_u F = -\tau^{-(L+1)}\sum_k w_k Y_k/(L-k+1)!$.
In the time $\sigma = t/\tau$, with $(\cdot)' = \dd(\cdot)/\dd\sigma$, the chain becomes
$z_j' = z_{j+1}$ for $j < L$ and $z_L' = \nu$, and the control equation becomes
$\nu' = -\sum_k w_k Y_k/(L-k+1)!$.
Since $z_j = z_1^{(j-1)}$ and $\nu = z_1^{(L)}$, the Laplace transform turns $Y_k$ into
$s^{k-1}e_{L-k+1}(s)$ times the transform of $z_1$, and the characteristic polynomial of the
resulting scalar equation for $z_1$ is $P$.
The change of variables is invertible for $\tau > 0$, and the time scaling multiplies the
eigenvalues by $1/\tau$.
\end{proof}

\begin{remark}[Weight design and update interval]\label{rem:chain}
For $L = 1$, $P = \sigma^2 + w_1\sigma + w_1$ is Hurwitz for every $w_1 > 0$.
For $L = 2$,
$P(\sigma) = \sigma^3 + \bigl(\frac{w_1}{4} + w_2\bigr)\sigma^2 + \bigl(\frac{w_1}{2} + w_2\bigr)\sigma
+ \frac{w_1}{2}$, which by the Routh--Hurwitz criterion is Hurwitz if and only if
$\bigl(\frac{w_1}{4} + w_2\bigr)\bigl(\frac{w_1}{2} + w_2\bigr) > \frac{w_1}{2}$.
Equal weights are admissible up to $L = 4$, which covers the chain from the rotor torques to the
position of a quadrotor, but not for $L = 5$, where $P$ with $w_k = 1$ has roots in the right
half-plane; longer chains call for a deliberate choice of the weights.
Since the closed-loop eigenvalues scale with $1/\tau$, the sampled-data loop \eqref{eq:update} is
stable only for $h/\tau$ below a threshold set by the roots of $P$.
For $L = 2$ with unit weights, it is stable if and only if $\tau > 1.526\,h$; a simulation that
advances the plant by the same forward Euler step would show the threshold $\tau > 1.408\,h$ instead.
Both are met by $\tau = mh$ with $m \ge 2$.
\end{remark}

\paragraph{Tracking error caused by the look-ahead.}
Even when the plant is on the reference, the look-ahead generally misses the target sample.
We call
\begin{equation}\label{eq:defect}
    \bdelta(t) := \bxref(t) - \hat\Phi_\tau\bigl(\bx_r(t),\bu_r(t)\bigr)
\end{equation}
the look-ahead defect: the amount by which the look-ahead misses the target when the plant is on the
reference and applies the reference input.
The defect is an object of the analysis only; the controller never evaluates it.

\begin{lemma}[Look-ahead defect]\label{lem:defect}
Let the reference stay in a compact set, let $\bx_r$ be three times and $\bu_r$ once continuously
differentiable, with bounded derivatives. If the target sample is the reference state at the time
$t + \theta$, $\bxref(t) = \bx_r(t+\theta)$ with $0 < \theta \le \bar\theta$, then, uniformly in $t$,
\begin{equation}\label{eq:defect-expansion}
    \bdelta(t) = (\theta - \tau)\,\dot\bx_r(t)
    + \tfrac12\bigl(\theta^2 - \tau^2\bigr)\ddot\bx_r(t)
    + \tfrac12\tau^2\,\partial_{\bu}\bv\bigl(\bx_r(t),\bu_r(t)\bigr)\,\dot\bu_r(t)
    + O(\theta^3 + \tau^3).
\end{equation}
For $\theta = \tau$, the defect is of second order only through the rate of the reference input, which
acts on the directly driven levels, and of third order otherwise.
\end{lemma}

\begin{proof}
Taylor expansion gives
$\bx_r(t+\theta) = \bx_r + \theta\dot\bx_r + \frac12\theta^2\ddot\bx_r + O(\theta^3)$, and
\eqref{eq:taylor-expansion} gives
$\hat\Phi_\tau(\bx_r,\bu_r) = \bx_r + \tau\bv + \frac12\tau^2\partial_{\bx}\bv\,\bv + O(\tau^3)$, with
$\bv$ and its derivatives evaluated at $(\bx_r,\bu_r)$.
Along the reference, $\dot\bx_r = \bv$ and $\ddot\bx_r = \partial_{\bx}\bv\,\bv + \partial_{\bu}\bv\,
\dot\bu_r$; subtracting yields \eqref{eq:defect-expansion}.
\end{proof}

The first term of \eqref{eq:defect-expansion} is of first order: unless the preview time matches the
time $\theta$ at which the reference reaches the target sample, the defect is dominated by a
mismatch of timing, which the closed loop absorbs as a time delay (\cref{ex:offset}).
With the next sample as target, $\theta$ is the sample interval, close to $h$, whereas the update
\eqref{eq:update} requires $\tau > 1.526\,h$ for a double integrator (\cref{rem:chain}).

A defect acts on the closed loop as a persistent disturbance, and so does the rate of the reference
input, which the control state has to follow.

\begin{theorem}[Ultimate bound on the tracking error]\label{thm:offset}
Let $\bxref = \hat\Phi_\tau(\bx_r,\bu_r) + \bdelta$ be a target with defect $\bdelta$, let
$\bxi = (\bx - \bx_r,\ \bu - \bu_r)$, and call nominal the error system obtained from \eqref{eq:hfc}
with $\bdelta = 0$ and without the term $-\dot\bu_r$; $\bxi = 0$ is an equilibrium of the nominal
system.
Suppose that this equilibrium is uniformly exponentially stable on a ball $B_r$, with a Lyapunov
function $V(t,\bxi)$ that satisfies $c_1\norm{\bxi}^2 \le V \le c_2\norm{\bxi}^2$,
$\dot V \le -c_3\norm{\bxi}^2$ along the nominal system, and $\norm{\partial_{\bxi}V} \le
c_4\norm{\bxi}$ on $B_r$. Let
\begin{equation}\label{eq:gbar}
    \bar g := \sup_{t,\ \bxi\in B_r}
        \norm[\big]{2D\,(\partial_{\bu}\hat\Phi_\tau)^\T W\bdelta(t)}
        + \sup_{t}\norm{\dot\bu_r(t)} .
\end{equation}
If $\bar g < \vartheta\,(c_3/c_4)\sqrt{c_1/c_2}\,r$ for some $\vartheta\in(0,1)$, then every solution of
\eqref{eq:hfc} with $\norm{\bxi(t_0)} < \sqrt{c_1/c_2}\,r$ satisfies, after a finite time,
\begin{equation}\label{eq:offset-bound}
    \norm{\bxi(t)} \le \frac{c_4}{c_3}\sqrt{\frac{c_2}{c_1}}\;\frac{\bar g}{\vartheta}.
\end{equation}
\end{theorem}

\begin{proof}
Since $\nabla_{\bu}F = 2(\partial_{\bu}\hat\Phi_\tau)^\T W\bigl(\hat\Phi_\tau(\bx,\bu) -
\hat\Phi_\tau(\bx_r,\bu_r) - \bdelta\bigr)$, the error dynamics are the nominal system plus the
perturbation
$\mathbf{g} = \bigl(\mathbf{0},\ 2D(\partial_{\bu}\hat\Phi_\tau)^\T W\bdelta - \dot\bu_r\bigr)$,
whose norm is bounded by $\bar g$ on $B_r$.
The claim is the ultimate-boundedness result for nonvanishing perturbations of exponentially stable
systems~\citep[Sect.~9.2]{khalil2002nonlinear}.
\end{proof}

By \cref{lem:defect}, the look-ahead leaves a tracking error driven by a disturbance of order
$\abs{\theta - \tau} + \tau^2\sup_t\norm{\dot\bu_r} + \tau^3$.
The constants $c_1,\dots,c_4$ and the coefficients $D$ depend on $\tau$, however, and the order of the
resulting offset depends on how they scale; for the integrator chain it can be computed exactly.
\TODO{一般系统中偏差阶数的严格证明（在 σ = t/τ 归一化时间下处理常数），见 notes/convergence.md}

\begin{example}[Delay of the double integrator]\label{ex:offset}
Take $L = 2$ in \cref{prop:chain} and a reference with constant acceleration $u_r = a_r$, whose
target sample at time $t$ is $\bx_r(t+\theta)$.
Substituting $x_1(t) = x_{1,r}(t - d) + c$, $x_2(t) = x_{2,r}(t - d)$ and $u = a_r$ into the
residuals of the cost shows that both vanish for all $t$ if and only if $d = \tau - \theta$ and
$c = 0$, whatever the weights, because the look-ahead is exact for a constant acceleration.
In steady state the plant therefore follows the reference with the time delay $\tau - \theta$ and
without offset; for $\theta = \tau$ it tracks the reference exactly.
The staged look-ahead of \cref{rem:staged} would leave the offset $-\frac12 a_r\tau^2$ instead.
\end{example}

\begin{remark}[Choice of the target sample]\label{rem:target}
With the next sample as target, $m = 1$, the time to the target is $\theta \approx h$, whereas the
update \eqref{eq:update} requires $\tau > 1.526\,h$ for a double integrator (\cref{rem:chain}); the
error \eqref{eq:error} is then dominated by the delay $\tau - h$, about $(\tau - h)\abs{\dot\bx_r}$ in
the units of the state.
Taking the $m$th next sample with $\tau = mh$, $m \ge 2$, matches the preview time to the time to the
target, so that the first-order term of \eqref{eq:defect-expansion} reduces to the deviation of the
sum of the next $m$ sample intervals from $mh$, while the update interval remains admissible
(\cref{rem:chain}).
What remains are this jitter of the sample intervals, the rate of the reference input, and terms of
third order.
\end{remark}

\subsection{Non-uniform sample intervals}\label{sec:nonuniform}

In the setting of \cref{sec:problem}, the target sample becomes due $m$ rounds after it is set, so
that $\tau = mh$ is the time the controller has to reach it, whatever the evolution times of the
reference between the samples.
If these times are markedly non-uniform, consuming one sample per round makes the samples
inconsistent with the clock of the controller: to reach every sample on time, the plant would have
to move faster or slower than the samples prescribe, and no choice of the preview time removes the
resulting error \eqref{eq:error}.
When the samples represent a trajectory $\bx_r(t)$ that is to be followed in time, their evolution
times can instead be estimated from the samples themselves.
If the rate of the lowest level does not depend on the input and stays away from zero, the time
between two samples is, up to terms of third order in the interval,
\begin{equation}\label{eq:interval}
    \widehat{\Delta t}_k = \frac{2\,\norm{\bx_{1,r,k+1} - \bx_{1,r,k}}}
        {\norm{\bv_1(\bx_{r,k})} + \norm{\bv_1(\bx_{r,k+1})}},
\end{equation}
the distance between the samples over their mean speed.
Summing these estimates gives the times $\hat t_j$ of the samples; at round $k$, the target is then the
first sample with $\hat t_j \ge t_k + mh$, the preview time is $\tau = \hat t_j - t_k$, and the
tracking error is the deviation from $\bx_r(t_k)$.
This selection needs neither time stamps nor a projection of the state onto the reference, and the
coefficients \eqref{eq:D} follow the preview time from round to round.
\TODO{仿真：按估计时间选目标时，相对 $\bx_r(t)$ 的误差与间隔抖动无关（notes Q39），数值在 Experiments 中报告}

\subsection{Design procedure}\label{sec:procedure}

\TODO{设计步骤（Algorithm 框）与参数选择：Taylor 预瞄阶数 L、$\tau = mh$（$m \ge 2$，满足\cref{rem:chain}的 τ > 1.526h）、权重（\cref{rem:chain}的 Hurwitz 条件）、前向 Euler 更新；自行车另见\cref{rem:bicycle-speed}}

\subsection{Example: swing-up and stabilization of a double inverted pendulum on a cart}
\label{sec:example-dipc}

\TODO{动力学模型（输入为作用在小车上的力）；两级层级，Taylor 预瞄阶数为 2；单一控制律完成从下垂位置起摆到收敛于倒立平衡点；控制律推导}

\subsection{Example: trajectory tracking with a kinematic bicycle model}
\label{sec:example-bicycle}

% Session 18: Taylor look-ahead, as in kinematics_flow_controller.hpp (refactored 2026-10-26);
% look-ahead, columns and rates checked against the code (notes/discussion.md, Q40).

\paragraph{Model.}
We describe the vehicle by the kinematic bicycle model~\citep{rajamani2012vehicle,kong2015kinematic},
with the reference point on the longitudinal axis at the distance $l_r$ ahead of the rear axle and
$l_f$ behind the front axle, for instance at the center of gravity; $l = l_f + l_r$ is the
wheelbase.
The state $\bx = (x,y,\psi,\delta,u)$ collects the position of the reference point in the world
frame, the heading, the steering angle of the front wheel, and the longitudinal velocity; the input
$\bu = (a,\omega)$ collects the longitudinal acceleration and the steering rate, which are bounded,
$\mathcal{U} = [a_{\min},a_{\max}]\times[\omega_{\min},\omega_{\max}]$; in this subsection, $\omega$
is the steering rate, not the speed of $Q$ in \cref{sec:properties}.
If neither wheel slips sideways, the yaw rate and the lateral velocity of the reference point follow
from $u$ and $\delta$ as $r = u\tan\delta/l$ and $v = l_r r$, and
\begin{equation}\label{eq:bicycle}
    \dot x = u\cos\psi - v\sin\psi, \quad
    \dot y = u\sin\psi + v\cos\psi, \quad
    \dot\psi = r, \quad
    \dot\delta = \omega, \quad
    \dot u = a,
\end{equation}
with $\abs{\delta} < \pi/2$.
The model is accurate as long as the lateral acceleration stays moderate, so that the tires barely
slip~\citep{kong2015kinematic,polack2017kinematic}.
Because the steering angle and the speed are states, the inputs are their rates, and the bounds of
$\mathcal{U}$ are the rate limits of the actuators.
The state falls into three levels, $\bx_1 = (x,y)$, $\bx_2 = \psi$ and $\bx_3 = (\delta,u)$: the
input drives the third level, the third level drives the heading, and the position is driven by the
heading and, through $u$ and $v$, directly by the third level.
The heading is a continuous function of time, for the plant as for the samples, so that heading
differences need no wrapping.

\paragraph{Look-ahead and cost.}
The Taylor look-ahead \eqref{eq:taylor} of \eqref{eq:bicycle} has order three.
Let $\mathbf{p} = (x,y)$, let $R(\psi)$ be the rotation by the heading and $J = R(\pi/2)$, and write
$\mathbf{b} = (u,v)$ for the velocity of the reference point in the body frame,
$\dot r = (a\tan\delta + u\,\omega\sec^2\delta)/l$ for the yaw acceleration, and
$\dot{\mathbf{b}} = (a, l_r\dot r)$ for the rate of $\mathbf{b}$.
Then
\begin{equation}\label{eq:bicycle-lookahead}
\begin{aligned}
    \delta^+ &= \delta + \tau\omega, \qquad u^+ = u + \tau a, \qquad
    \psi^+ = \psi + \tau r + \tfrac12\tau^2\dot r, \\
    \mathbf{p}^+ &= \mathbf{p} + R(\psi)\Bigl(\tau\mathbf{b}
    + \tfrac12\tau^2\bigl(\dot{\mathbf{b}} + rJ\mathbf{b}\bigr)
    + \tfrac16\tau^3\bigl(2rJ\dot{\mathbf{b}} + \dot r J\mathbf{b} - r^2\mathbf{b}\bigr)\Bigr).
\end{aligned}
\end{equation}
The terms up to second order are the derivatives of the flow, the velocity of the reference point
turning with the heading; at third order, the second derivative $\ddot r$ of the yaw rate, which is
quadratic in $(a,\omega)$, is omitted in the position and in the heading.
With the steering angle held, $\omega = 0$, the look-ahead is therefore exact through $\tau^3$, and
along a straight line under constant acceleration it is exact.
Every component of \eqref{eq:bicycle-lookahead} is affine in $(a,\omega)$.

The target $\bxref = (x_{\mathrm{ref}},y_{\mathrm{ref}},\psi_{\mathrm{ref}},\delta_{\mathrm{ref}},
u_{\mathrm{ref}})$ is the $m$th next sample \eqref{eq:xref}, with $\tau = mh$, and the residuals are
weighted as
\begin{equation}\label{eq:bicycle-cost}
    F = w_p\,\frac{(x^+ - x_{\mathrm{ref}})^2 + (y^+ - y_{\mathrm{ref}})^2}{l^2}
    + w_\psi\,(\psi^+ - \psi_{\mathrm{ref}})^2 + w_\delta\,(\delta^+ - \delta_{\mathrm{ref}})^2
    + w_u\,\frac{\tau^2(u^+ - u_{\mathrm{ref}})^2}{l^2},
\end{equation}
that is, $\bar\Gamma = \diag(l^{-1},l^{-1},1,1,\tau/l)$ in \eqref{eq:cost}: positions are measured
in wheelbases, and the speed by the distance it covers in the preview time, again in wheelbases.
Equal weights on $x$ and $y$ keep $F$ invariant under rotations of the world frame.
Since $\partial_a\bv$ and $\partial_\omega\bv$ are the unit vectors of the components $u$ and
$\delta$, the coefficients \eqref{eq:D} are $d_a = l^2/(2\tau^5)$ and $d_\omega = 1/(2\tau^3)$.

Because the look-ahead is affine in the input,
$\nabla_{\bu}F = 2(\partial_{\bu}\hat\Phi_\tau)^\T W(\hat\Phi_\tau - \bxref)$ with
$W = \diag(w_p/l^2, w_p/l^2, w_\psi, w_\delta, w_u\tau^2/l^2)$, and the columns of
$\partial_{\bu}\hat\Phi_\tau$, with the components ordered as $(x,y,\psi,\delta,u)$, are
\begin{equation}\label{eq:bicycle-columns}
\begin{aligned}
    \partial_a\hat\Phi_\tau &= \tau^2\Bigl(
        \frac12 R(\psi)\begin{bmatrix} 1 - \tau v\tan\delta/l \\ (l_r + \tau u)\tan\delta/l \end{bmatrix},\
        \frac{\tan\delta}{2l},\ 0,\ \frac{1}{\tau}\Bigr), \\
    \partial_\omega\hat\Phi_\tau &= \tau^2\Bigl(
        \frac{u\sec^2\delta}{l}\,R(\psi)\begin{bmatrix} -\tau v/2 \\ l_r/2 + \tau u/6 \end{bmatrix},\
        \frac{u\sec^2\delta}{2l},\ \frac{1}{\tau},\ 0\Bigr).
\end{aligned}
\end{equation}
The \method{} law \eqref{eq:hfc} becomes
\begin{equation}\label{eq:bicycle-law}
    \dot a = -\frac{l^2}{\tau^5}\,(\partial_a\hat\Phi_\tau)^\T W\bigl(\hat\Phi_\tau - \bxref\bigr),
    \qquad
    \dot\omega = -\frac{1}{\tau^3}\,(\partial_\omega\hat\Phi_\tau)^\T W\bigl(\hat\Phi_\tau - \bxref\bigr),
\end{equation}
which the controller integrates by the update \eqref{eq:update}.
A round of the controller thus takes a fixed number of arithmetic operations and a few
trigonometric functions; it involves neither a matrix inversion nor an iteration.

\paragraph{How the errors reach the input.}
With the forward Euler look-ahead, only the residuals of $\delta$ and $u$ would enter
$\nabla_{\bu}F$, and the controller would follow the steering angle and the speed of the samples
without any feedback from heading and position (\cref{prop:sg}).
By \eqref{eq:bicycle-columns}, the Taylor look-ahead lets the heading error enter at order $\tau^2$,
through the dependence of the yaw rate on $\delta$ and $u$, and the position error as well, through
the dependence of the velocity of the reference point on $u$ and, sideways, on $\delta$.
The lateral entry of $\partial_\omega\hat\Phi_\tau$ in the body frame,
$\tau^2 u\sec^2\delta\,(l_r/2 + \tau u/6)/l$, combines two effects: the lateral velocity of the
reference point, which builds up over the preview time with the lever arm $l_r/2$, and the rotation of
the velocity by the heading change, with the lever arm $\tau u/6$.
The second effect, of third order, dominates whenever $\tau u > 3l_r$, and it is the only one when the
reference point is at the rear axle, $l_r = 0$; a look-ahead truncated at second order would leave
the lateral position of such a vehicle without any influence on the steering rate.

\paragraph{Straight references.}
About a straight reference traversed at constant speed, the longitudinal motion reduces to the
integrator chain of \cref{prop:chain}, and the lateral motion to a loop of fourth order.

\begin{proposition}[Bicycle on a straight reference]\label{prop:bicycle}
Linearize the closed loop \eqref{eq:bicycle}, \eqref{eq:bicycle-law} about a straight reference
traversed at the constant speed $U > 0$, whose target at the time $t$ is its state at $t + \tau$, and
let $\kappa = U\tau/l$ be the preview distance in wheelbases.
\begin{enumerate}
    \item The linearization splits into a longitudinal part, in the along-track position, $u$ and $a$,
    and a lateral part, in the cross-track position, $\psi$, $\delta$ and $\omega$.
    \item The longitudinal part is the loop of \cref{prop:chain} with $L = 2$, $\ell_0 = l$,
    $w_1 = w_p$ and $w_2 = w_u$.
    \item The characteristic polynomial of the lateral part is $\tau^{-4}P(\tau s)$ with
    \begin{equation}\label{eq:bicycle-lateral}
        P(\sigma) = \sigma^4 + a_3\sigma^3 + a_2\sigma^2 + a_1\sigma + a_0,
    \end{equation}
    where, with $\beta = \kappa\,(\kappa/6 + l_r/(2l))$,
    \[
    \begin{aligned}
        a_0 &= w_p\beta\kappa^2, &
        a_1 &= w_p\beta\kappa\,(\kappa + l_r/l) + \tfrac12 w_\psi\kappa^2, \\
        a_2 &= w_p\beta\kappa\,(\kappa/2 + l_r/l) + \tfrac12 w_\psi\kappa^2 + w_\delta, &
        a_3 &= w_p\beta^2 + \tfrac14 w_\psi\kappa^2 + w_\delta .
    \end{aligned}
    \]
    \item For unit weights, both parts are exponentially stable for every $\kappa > 0$ and every
    $l_r\in[0,l]$, so that the reference is locally exponentially stable.
\end{enumerate}
\end{proposition}

\begin{proof}
We place the reference on the $x$-axis, where $\psi = \delta = v = 0$, $u = U$ and $a = \omega = 0$,
and denote the deviations from it by the same letters.
The reference input lies in the interior of $\mathcal{U}$, so the clamping is inactive near it.

(a) The lateral deviations enter the longitudinal components of \eqref{eq:bicycle} and
\eqref{eq:bicycle-lookahead} only in products of deviations, and vice versa.
In \eqref{eq:bicycle-columns}, the entries that couple the two parts carry the factor $\tan\delta$
or $v$, so they are of first order, and their products with the residuals drop out of the
linearization of \eqref{eq:bicycle-law}.

(b) The longitudinal part reads $\dot x = u$, $\dot u = a$, $x^+ = x + \tau u + \frac12\tau^2 a$ and
$u^+ = u + \tau a$, with the cost $w_p(x^+ - x_{\mathrm{ref}})^2/l^2 +
w_u\tau^2(u^+ - u_{\mathrm{ref}})^2/l^2$ and $d_a = l^2/(2\tau^5)$.
In the deviations from the reference, whose target is exactly its own look-ahead, this is the loop
of \cref{prop:chain}.

(c) Let $\eta = y/l$, $\nu = \tau\omega$, and $(\cdot)' = \dd(\cdot)/\dd\sigma$ with $\sigma = t/\tau$.
To first order, $\dot r = U\omega/l$ and $v = l_r U\delta/l$, so the lateral part reads
$\eta' = \kappa(\psi + (l_r/l)\delta)$, $\psi' = \kappa\delta$ and $\delta' = \nu$, and its look-ahead
\[
    \frac{y^+}{l} = \eta + \kappa\psi + \kappa\Bigl(\frac{\kappa}{2} + \frac{l_r}{l}\Bigr)\delta
    + \beta\nu,
    \qquad \psi^+ = \psi + \kappa\delta + \tfrac12\kappa\nu, \qquad \delta^+ = \delta + \nu .
\]
Since $\partial_\omega = \tau\partial_\nu$, the control equation is
$\nu' = \tau^2\dot\omega = -\frac12\partial_\nu F
= -w_p\beta\,y^+/l - \frac12 w_\psi\kappa\psi^+ - w_\delta\delta^+$.
In the Laplace domain, $\nu = s\delta$, $\psi = \kappa\delta/s$ and
$\eta = \kappa(\kappa + (l_r/l)s)\delta/s^2$; substituting these into the control equation and
multiplying by $s^2$ gives $P(s)\delta = 0$.
The scaling of time multiplies the eigenvalues by $1/\tau$.

(d) For unit weights, the longitudinal part satisfies
$\bigl(\frac{w_1}{4} + w_2\bigr)\bigl(\frac{w_1}{2} + w_2\bigr) = \frac{15}{8} > \frac{w_1}{2}$
(\cref{rem:chain}).
For the lateral part, the coefficients $a_0,\dots,a_3$ are positive, and the Hurwitz determinants
$a_3a_2 - a_1$ and $a_1(a_3a_2 - a_1) - a_3^2a_0$ are polynomials in $\kappa$ and $l_r/l$;
quantifier elimination over the reals by cylindrical algebraic decomposition~\citep{collins1975quantifier}
shows that both are positive for all $\kappa > 0$ and $l_r/l\in[0,1]$.
Hence $P$ is Hurwitz, and Lyapunov's indirect method gives the local exponential stability.
\end{proof}

\begin{remark}[Speed range of the sampled-data loop]\label{rem:bicycle-speed}
With $\tau = mh$ and $m \ge 2$, the longitudinal part meets the condition $\tau > 1.526\,h$ of
\cref{rem:chain}.
The lateral part bounds the preview distance instead.
For unit weights and $l_r = l/2$, the sampled-data loop \eqref{eq:update}, linearized about a straight
reference, has spectral radius below one for $\kappa < 0.92$, $1.21$ and $3.37$ with $m = 2$, $3$ and
$4$; for $l_r = 0$, the bounds are $1.15$, $1.49$ and $3.94$.
These bounds depend on $\kappa$, $l_r/l$ and $m$ only; the loop can regain stability beyond them, but
we do not rely on it.
In practice, the preview distance $U\tau$ should therefore stay below about one wheelbase; at higher
speeds, the update interval $h$ has to be shortened or the weight $w_p$ lowered.
\end{remark}

\paragraph{Tracking error.}
By \cref{lem:defect} with $\theta = \tau$, the defect is
$\frac12\tau^2\partial_{\bu}\bv\,\dot\bu_r + O(\tau^3)$: it is of second order only in the steering
angle and the speed, $\frac12\tau^2\dot\omega_r$ and $\frac12\tau^2\dot a_r$, and of third order
otherwise.
On references with a held steering angle and a constant acceleration, which include straight lines and
circles traversed at constant speed, the look-ahead is exact through $\tau^3$ and the defect is of
fourth order, so that the steady offsets essentially vanish, whereas the staged look-ahead of
\cref{rem:staged} leaves the vehicle about $\frac12\tau^2$ times the centripetal acceleration outside
a curve.
With $\tau = mh$, the error \eqref{eq:error} is then dominated by the jitter of the sample intervals,
which enters at first order (\cref{rem:target}).
\TODO{稳态偏差的仿真数值（圆周、匀加速直线）在 Experiments 中报告}

\paragraph{Conditions on the reference.}
The samples carry all five components of the state; the controller uses no input of the reference.
For the analysis, the samples are taken from a solution of \eqref{eq:bicycle} with continuously
differentiable inputs in the interior of $\mathcal{U}$, $\abs{\delta_r} < \pi/2$, and a speed bounded
away from zero.
The last condition is essential: as $U \to 0$, $\kappa \to 0$ and the lateral feedback in
\eqref{eq:bicycle-lateral} vanishes with $a_0$ and $a_1$.
This is not a shortcoming of the particular law: by Brockett's necessary
condition~\citep{brockett1983asymptotic}, no continuous time-invariant feedback, static or dynamic,
asymptotically stabilizes a parked pose of the vehicle, since no state near rest and no input
produce a sideways velocity.

\subsection{Example: trajectory tracking of a quadrotor unmanned aerial vehicle}
\label{sec:example-quadrotor}

\TODO{模型（平移 + 姿态，姿态在 SO(3) 上完整描述）；四级层级，Taylor 预瞄阶数为 4，姿态部分用 SO(3) 上的展开（Magnus 型），输入在 u = 0 处线性化；参考轨迹的条件；控制律推导}


%----------------------------------------------------------------------------------------
%	SECTION 3: EXPERIMENTS
%----------------------------------------------------------------------------------------

\section{Experiments}\label{sec:experiments}

% Second stage: to be written once the simulation and experimental results are available
% (plan in notes/outline.md, Section 3).
\TODO{第二阶段：仿真与实验结果齐备后撰写}


%----------------------------------------------------------------------------------------
%	SECTION 4: CONCLUSION
%----------------------------------------------------------------------------------------

\section{Conclusion}\label{sec:conclusion}

% Second stage: summary of the method, scope of generalization, possible variants
% (plan in notes/outline.md, Section 4).
\TODO{第二阶段：实验完成后撰写}


%----------------------------------------------------------------------------------------
%	ACKNOWLEDGEMENTS
%----------------------------------------------------------------------------------------

\section*{Acknowledgements}

\TODO{资助信息与致谢}


%----------------------------------------------------------------------------------------
%	DECLARATIONS
%----------------------------------------------------------------------------------------

\section*{Declarations}

% Springer "Statements and Declarations"; confirm the exact headings required by the
% Nonlinear Dynamics submission guidelines before submission.

\paragraph{Funding.}
\TODO{资助信息}

\paragraph{Competing interests.}
\TODO{利益冲突声明}

\paragraph{Data availability.}
\TODO{数据可用性声明}

\paragraph{Code availability.}
\TODO{代码可用性声明（是否公开仿真与控制代码）}


%----------------------------------------------------------------------------------------
%	BIBLIOGRAPHY
%----------------------------------------------------------------------------------------

\printbibliography

\end{document}
